\clearpage\thispagestyle{empty}
\pdfbookmark[0]{Quelle zum Brief an von Lindenau vom 27. Juni 1808}{lindenau1808june-sources}
\begin{center}\textbf{Quelle zum Brief an von Lindenau vom 27. Juni 1808}\end{center}
\noindent Diese KI-Übertragung ins heutige Deutsch folgt dem Brief von Carl Friedrich Gauß an Bernhard August von Lindenau vom 27. Juni 1808 in \href{https://gdz.sub.uni-goettingen.de/id/PPN236007467}{Gauß’ \textit{Werken}, Bd.~VI} (Göttingen, 1874), S.~301--303. Der Brief erschien zuerst als Beitrag~IX, \textit{Auszug aus einem Schreiben des Hrn. Professors Gauß}, im Juliheft 1808 der \textit{Monatlichen Correspondenz}, Bd.~XVIII, \href{https://www.digitale-sammlungen.de/en/view/bsb10538609?page=95}{S.~83}, \href{https://www.digitale-sammlungen.de/en/view/bsb10538609?page=96}{84}, \href{https://www.digitale-sammlungen.de/en/view/bsb10538609?page=97}{85} und \href{https://www.digitale-sammlungen.de/en/view/bsb10538609?page=98}{86}. Das digitalisierte Exemplar stammt aus der Bayerischen Staatsbibliothek.
\par\medskip
Im Erstdruck misst Hardings Dollond-Instrument $3\tfrac12$ Fuß, in der Werkausgabe $3\tfrac14$ Fuß. Die Tafeln folgen der Werkausgabe. Im Brief steht das Aphel der Vesta, im gesonderten Bericht vom 4. Juli das Perihel. Auch die Schreibweise \textit{Zeifel} und das Komma nach \textit{lasse} folgen der Werkausgabe.
\par\medskip
Gauß behandelt die Bahnbestimmung aus vier Beobachtungen in der \textit{Theoria motus}, Buch~II, Abschnitt~II. In Charles Henry Davis’ englischer Übersetzung von 1857 stehen die Aufgabenstellung in \href{https://www2.karlin.mff.cuni.cz/~tuma/Aplikace21/theoryofmotionof00gausuoft.pdf#page=258}{Art.~164} und das \href{https://www2.karlin.mff.cuni.cz/~tuma/Aplikace21/theoryofmotionof00gausuoft.pdf#page=267}{Vesta-Beispiel in Art.~171}, S.~243--248. Das Buch enthält seine eigene Rechnung und eigene Tafelangaben. KI-Übertragung und KI-Satz.
\clearpage
